\documentclass[10pt,a4paper]{amsart}
\usepackage{amsmath,amssymb,amsthm,graphicx}
% Retained minimal preamble for the Stacks Project chapter more-algebra.tex (More on Algebra).
% The chapter file pinned-original-source.tex carries only \input{preamble} on its first line, so no
% preamble text lives inside the chapter itself.  The Stacks Project uses one single shared preamble for
% every chapter; the snapshot stages it identically as sources/stacks-chow/preamble.tex and
% sources/stacks-etale-cohomology/preamble.tex.  The block below is copied mechanically, line for line,
% from sources/stacks-chow/preamble.tex lines 14-18 (the Xy-pic setup) and 154-217 (the native theorem
% declarations together with the mathematical macros).  Nothing was added, removed or redefined; the
% wrapper class and the amsmath/amssymb/amsthm/graphicx base packages are supplied by the assembler.
% Chapter: more-algebra.tex -- More on Algebra.
% For commutative diagrams we use Xy-pic
\usepackage[all]{xy}

% We use 2cell for 2-commutative diagrams.
\xyoption{2cell}

% Theorem environments.
%
\theoremstyle{plain}
\newtheorem{theorem}[subsection]{Theorem}
\newtheorem{proposition}[subsection]{Proposition}
\newtheorem{lemma}[subsection]{Lemma}

\theoremstyle{definition}
\newtheorem{definition}[subsection]{Definition}
\newtheorem{example}[subsection]{Example}
\newtheorem{exercise}[subsection]{Exercise}
\newtheorem{situation}[subsection]{Situation}

\theoremstyle{remark}
\newtheorem{remark}[subsection]{Remark}
\newtheorem{remarks}[subsection]{Remarks}

\numberwithin{equation}{subsection}

% Macros
%
\def\lim{\mathop{\mathrm{lim}}\nolimits}
\def\colim{\mathop{\mathrm{colim}}\nolimits}
\def\Spec{\mathop{\mathrm{Spec}}}
\def\Hom{\mathop{\mathrm{Hom}}\nolimits}
\def\Ext{\mathop{\mathrm{Ext}}\nolimits}
\def\SheafHom{\mathop{\mathcal{H}\!\mathit{om}}\nolimits}
\def\SheafExt{\mathop{\mathcal{E}\!\mathit{xt}}\nolimits}
\def\Sch{\mathit{Sch}}
\def\Mor{\mathop{\mathrm{Mor}}\nolimits}
\def\Ob{\mathop{\mathrm{Ob}}\nolimits}
\def\Sh{\mathop{\mathit{Sh}}\nolimits}
\def\NL{\mathop{N\!L}\nolimits}
\def\CH{\mathop{\mathrm{CH}}\nolimits}
\def\proetale{{pro\text{-}\acute{e}tale}}
\def\etale{{\acute{e}tale}}
\def\QCoh{\mathit{QCoh}}
\def\Ker{\mathop{\mathrm{Ker}}}
\def\Im{\mathop{\mathrm{Im}}}
\def\Coker{\mathop{\mathrm{Coker}}}
\def\Coim{\mathop{\mathrm{Coim}}}

% Boxtimes
%
\DeclareMathSymbol{\boxtimes}{\mathbin}{AMSa}{"02}

%
% Macros for moduli stacks/spaces
%
\def\QCohstack{\mathcal{QC}\!\mathit{oh}}
\def\Cohstack{\mathcal{C}\!\mathit{oh}}
\def\Spacesstack{\mathcal{S}\!\mathit{paces}}
\def\Quotfunctor{\mathrm{Quot}}
\def\Hilbfunctor{\mathrm{Hilb}}
\def\Curvesstack{\mathcal{C}\!\mathit{urves}}
\def\Polarizedstack{\mathcal{P}\!\mathit{olarized}}
\def\Complexesstack{\mathcal{C}\!\mathit{omplexes}}
% \Pic is the operator that assigns to X its picard group, usage \Pic(X)
% \Picardstack_{X/B} denotes the Picard stack of X over B
% \Picardfunctor_{X/B} denotes the Picard functor of X over B
\def\Pic{\mathop{\mathrm{Pic}}\nolimits}
\def\Picardstack{\mathcal{P}\!\mathit{ic}}
\def\Picardfunctor{\mathrm{Pic}}
\def\Deformationcategory{\mathcal{D}\!\mathit{ef}}
\title{Base change of the naive cotangent complex of a presentation is an isomorphism on H\textasciicircum{}0 and onto on H\textasciicircum{}\{-1\}}
\author{The Stacks Project Authors}
\date{stacks-project:more-algebra.tex:local-snapshot-20261012; GNU FDL 1.2 or later}
\begin{document}
\maketitle
\setcounter{subsection}{1}
\noindent\textit{Editorial note.} All delivered mathematical text is the author's own, copied from the pinned Stacks Project chapter file more-algebra.tex; no statement or proof step was rewritten, reordered or removed. The counted statement is the labeled lemma environment at source lines 22770-22786 (canonical tag 0FUZ) and the counted proof is the proof environment at source lines 22788-22807, which immediately follows it and is its own proof.\par\medskip
\noindent\textit{Editorial note.} The macro NL of the naive cotangent complex is supplied by the retained preamble, which copies the shared Stacks preamble verbatim; the displayed square of the proof is the author's own Xy-pic markup.\par\medskip
\noindent\textit{Editorial note.} The proof transfers the presentation-level statement to the canonical complexes by using that the two are homotopic, a fact the author invokes without a label; the transfer is the last step and is preserved verbatim, while the substantive comparison is carried out for a polynomial presentation inside the counted span.\par\medskip
\noindent\textit{Editorial note.} The chapter file pinned-original-source.tex carries only an input of the shared preamble on its first line, so no preamble text lives inside the chapter. The retained file sources/stacks-more-algebra/retained-preamble.tex is a mechanical, line-for-line copy of the project's single shared preamble as staged in the snapshot (Xy-pic setup plus the native theorem declarations and mathematical macros). Nothing was added, removed or redefined, and the excerpt serves only as preamble, never as a span coordinate.\par\medskip
\noindent\textit{Editorial note.} The source URL is the source-level Stacks Project URL, as required when more than one card is produced from one source; the per-claim tag is carried by claim\_key and by this note. The counted statement and proof cite no bibliography entry, so this card needs no bibliography.\par\medskip
\noindent\textit{Allowed standard prerequisites (editorial).}\begin{itemize}
\item The naive cotangent complex of a presentation P -> S with kernel I is the two-term complex I/I\textasciicircum{}2 -> S tensor over P for the module of differentials, displayed in degrees -1 and 0. used at proof lines 22789-22800, where I and I' are introduced and the displayed square is identified with NL(alpha) tensor S' on the left and NL(alpha') on the right. Fixes the explicit two-term complexes whose base change map has to be analysed.
\item Tensor product is right exact, so the natural map I tensor\_R R' -> I' is surjective when I' is defined as the kernel of P tensor R' -> S tensor R'. used at proof lines 22801-22803. Makes the bottom horizontal arrow of the displayed square surjective, which is the assertion about H\textasciicircum{}\{-1\}.
\item The naive cotangent complex of a ring map is independent of the chosen presentation up to homotopy, so the two complexes in the second displayed map are homotopic to the complexes of a fixed presentation. invoked at proof lines 22804-22806, where the statement for NL\_\{S/R\} is deduced because the complexes are homotopic to NL(alpha) tensor S' and to NL(alpha'). Transfers the presentation-level statement to the canonical naive cotangent complex of the ring map, which is the second assertion of the lemma.
\end{itemize}\medskip
\begin{lemma}
\label{lemma-base-change-NL}
Let $R \to S$ and $R \to R'$ be ring maps.
Let $\alpha : P \to S$ be a presentation of $S$ over $R$.
Then $\alpha' : P \otimes_R R' \to S \otimes_R R'$ is a
presentation of $S' = S \otimes_R R'$ over $R'$.
The canonical map
$$
NL(\alpha) \otimes_S S' \to \NL(\alpha')
$$
is an isomorphism on $H^0$ and surjective on $H^{-1}$. In particular,
the canonical map
$$
\NL_{S/R} \otimes_S S' \to \NL_{S'/R'}
$$
is an isomorphism on $H^0$ and surjective on $H^{-1}$.
\end{lemma}

\begin{proof}
Denote $I = \Ker(P \to S)$. Denote $P' = P \otimes_R R'$ and
$I' = \Ker(P' \to S')$. Suppose $P$ is a polynomial algebra
on $x_j$ for $j \in J$. The map displayed in the lemma becomes
$$
\xymatrix{
\bigoplus_{j \in J} S' \text{d}x_j \ar[r] &
\bigoplus_{j \in J} S' \text{d}x_j \\
I/I^2 \otimes_S S' \ar[r] \ar[u] &
I'/(I')^2 \ar[u]
}
$$
where the left column is $\NL(\alpha) \otimes_S S'$ and the right
column is $\NL(\alpha')$. By right exactness of tensor product we
see that $I \otimes_R R' \to I'$ is surjective. Hence the bottom arrow
is a surjection. This proves the first statement of the lemma.
The statement for $\NL_{S/R} \otimes_S S' \to \NL_{S'/R'}$
follows as these complexes are homotopic to $\NL(\alpha) \otimes_S S'$ and
$\NL(\alpha')$.
\end{proof}

\end{document}
